% BEGIN SELECTION RISK AUDIT
\subsubsection{Return-Oriented versus Defensive Selection}\label{sec:selectionrisk}
We separately compare the supplied momentum--profitability--issuance selection (50 stocks) with the low-risk--coskewness defensive selection (30 stocks). Both use February 2003--December 2020, observed CRSP total returns, drift-adjusted transaction costs of 10 basis points per dollar traded, and terminal proceeds held in cash. Constituents are retained exactly; no new size filter or missing-return imputation is applied. The two allocation benchmarks are equal weighting and recorded market-capitalization weights, with the common rule $u=2/N$ giving limits of 4.00\% and 6.67\%, respectively. Reconstructed proportional weights agree with both supplied files within $5\times10^{-7}$, and the defensive monthly net returns reproduce the article baseline within $10^{-12}$.

\begin{table*}[!htbp]\centering
\caption{Selection comparison under the article's return and cost conventions. CAGR, annualized volatility, and maximum drawdown (MDD) are percentages; Sharpe is dimensionless. Selection and portfolio size change together. These are allocation benchmarks, not a re-estimation of all optimizers on the new universe.}\label{tab:selectionrisk}
\small\begin{tabular}{llrrrr}\toprule
Selection ($N$) & Allocation & CAGR & Volatility & Sharpe & MDD\\\midrule
Defensive (30) & Equal weight & 13.52 & 11.18 & 1.087 & -30.77 \\
Defensive (30) & Market-cap weight & 14.05 & 11.39 & 1.110 & -29.00 \\
Return-oriented (50) & Equal weight & 16.60 & 20.37 & 0.801 & -49.18 \\
Return-oriented (50) & Market-cap weight & 16.48 & 18.45 & 0.858 & -48.05 \\
\bottomrule\end{tabular}\end{table*}

Under market-capitalization weighting, CAGR rises by 2.43 percentage points, but volatility rises by 7.06 points and Sharpe falls from 1.110 to 0.858. Maximum drawdown becomes 19.04 percentage points more severe. Equal weighting gives the same qualitative result. Thus the observed growth gain comes with materially greater realized risk; the new selection does not establish a return improvement at comparable risk.

For market-capitalization weights, a paired six-month circular block bootstrap (2,000 draws) gives pointwise 95\% intervals of $[-4.46,9.58]$ percentage points for the CAGR difference, $[5.18,8.91]$ points for the volatility difference, and $[-0.628,0.073]$ for the Sharpe difference. All differences are return-oriented minus defensive. The volatility increase is clear, whereas the CAGR interval includes zero. The paired mean-return HAC(4) test has $p=0.312$, or $p=0.401$ after Holm adjustment across the two weighting rules. A Sharpe interval containing zero is not an equivalence test or evidence that risk is unchanged.

In 2012--2020, equal-weight CAGR is 12.85\% for the return-oriented selection versus 12.98\% for the defensive selection. Market-capitalization weighting retains a CAGR increase (14.99\% versus 13.56\%), but volatility remains higher (15.86\% versus 11.40\%) and Sharpe lower (0.924 versus 1.123). The comparison is conditional on historical selections: different stock counts, associated position limits, and prior factor searches prevent attributing the difference solely to the factors. Optimization may still improve risk-adjusted performance within the new selection, but that is a separate comparison against its own equal-weight baseline.
% END SELECTION RISK AUDIT

